\chapter*{General Introduction}
\addcontentsline{toc}{chapter}{General Introduction}
\chaptermark{General Introduction}
\label{chap:general_introduction}

\section*{Project Background}
\addcontentsline{toc}{section}{Project Background}

Colorectal cancer is still a major public health problem worldwide \cite{sung2021globocan}.
One of the best ways to reduce that risk is to remove suspicious polyps early, which is
why \textbf{polypectomy} matters so much \cite{winawer1993polypectomy}. Even though the
procedure is common, it is not simple: the final result depends on the operator, the
quality of the view, the shape of the lesion, and the pressure of making good decisions
in real time.

Even with good screening programs, tandem colonoscopy studies still report meaningful
\textbf{adenoma} (a type of polyp that can become cancerous) miss rates
\cite{zhao2019adenomaMiss}. Also, quality measures like \textbf{adenoma detection rate}
are linked with future colorectal cancer risk \cite{corley2014adr}.

This work studies how intelligent methods can support \textbf{polypectomy} by
combining computer vision and \textbf{reinforcement learning} (RL, learning by trial and
error using rewards). The main idea is to help an agent read the scene, understand what
is happening, and choose better actions in a simulated setting. The goal here is to
study how perception and action work together in a sequence of decisions, not to build a
clinical detector.

\section*{Problem}
\addcontentsline{toc}{section}{Problem}
\vspace{-0.5em}
The core problem is to design a \textbf{learning-based control and perception framework}
for \textbf{polypectomy-like tasks} that works under realistic constraints:
\textbf{partial and noisy visual observations}, \textbf{delayed or limited instrument control},
and strict \textbf{safety requirements}. Concretely, the main technical challenges are to:
interpret noisy imagery reliably, decide actions that avoid tissue damage, time interventions
under uncertainty, and remain robust when the visual scene changes or tools occlude
important features.

Concretely, this work addresses the following specific questions:
\begin{itemize}
\item \textbf{Perception → Action:} How can an agent learn to use visual observations
	(single frames, short temporal stacks, or depth cues) to select \textbf{safe and effective}
	actions during a polypectomy-like task?
\item \textbf{Learning & Transfer:} How can training be organised so the model improves
	in \textbf{simulation} while enforcing \textbf{safety limits} and conservative procedures
	for later transfer to real hardware?
\item \textbf{Reward & Representation:} How should the \textbf{reward function} and
	\textbf{observation encoding} be designed so the agent balances \textbf{tissue preservation},
	precision of resection, and task efficiency?
\item \textbf{Robustness:} How to make policies resilient to \textbf{domain shift} such as
	camera viewpoint changes, lighting variability, specular highlights, and instrument occlusion?
\item \textbf{Safety Mechanisms:} What \textbf{safety checks}, motion limits, and fallback
	strategies are required to prevent harmful actions during transfer from simulation to
	a physical endoscope?
\item \textbf{Evaluation:} Beyond simple success/failure, which metrics (\textbf{stability},
	\textbf{precision}, \textbf{efficiency}, and \textbf{clinician acceptability}) best capture
	practical usefulness?
\end{itemize}

This problem is motivated by the project theme: build a \textbf{learning-based system}
for sequential decision-making in \textbf{simulated polypectomy tasks}. The aim is to
explore technical solutions that prioritise safety, interpretability, and robust
behaviour under realistic visual and control constraints.

\section*{Objectives}
\addcontentsline{toc}{section}{Objectives}
The objectives of this work are:
\begin{itemize}
\item to review the clinical and technical background of \textbf{polypectomy} and the main
intelligent methods linked to it;
\item to study how machine learning, deep learning, and \textbf{reinforcement learning}
can help medical decision-making;
\item to design a learning-based framework that fits sequential decision-making in
polyp removal;
\item to identify the evaluation criteria that matter most for performance, stability,
and practical usefulness.
\end{itemize}

\section*{Proposed Approach}
\addcontentsline{toc}{section}{Proposed Approach}
In this work, we propose a learning-based framework that combines simple visual
processing with sequential decision learning. The main objectives are to build a safe
simulation environment, train agents that can act reliably in that environment, and
evaluate how well learned policies generalize across cases.

Main solutions and steps:
\begin{itemize}
\item \textbf{Simulation and data:} build a Unity-based simulated environment with
predefined targets and controls for the endoscope and tools.
\item \textbf{Learning pipeline:} use convolutional networks (CNNs) for visual encoding
and reinforcement learning (PPO) for sequential control, with careful reward shaping.
\item \textbf{Safety and limits:} design safety checks, motion limits, and conservative
transfer steps for any prototype testing.
\item \textbf{Evaluation:} define clear metrics (success rate, stability, efficiency) and
compare learned behavior to baseline scripted or expert demonstrations.
\end{itemize}

These solutions target the problem by combining perception, action, and safe practice
within a controlled simulation.

\section*{Document Plan}
\addcontentsline{toc}{section}{Document Plan}
The document is arranged in two main parts that follow from the repository's scope:

\begin{itemize}
\item \textbf{Chapter 1 --- State of the Art:} a concise review that situates the
	clinical problem (colorectal cancer and polypectomy), summarises key clinical
	metrics and risks, and surveys technical approaches from computer vision and
	reinforcement learning that are relevant to the project. This chapter highlights
	practical gaps and limitations in current methods that motivate the system design.

\item \textbf{Chapter 2 --- Contributions:} an account of the implemented system and
	the original contributions. In short, this chapter presents an attempt on a
		extbf{two-phase endoscope reinforcement learning system} that combines a
		extbf{simulation environment}, a \textbf{perception module} (visual encoder / CNN),
	and a \textbf{two-level control strategy} (coarse navigation + fine manipulation)
	trained with reinforcement learning. The chapter documents the design choices,
	safety measures, training setup, and evaluation metrics used to assess stability,
	precision, and practical usefulness.
\end{itemize}

Overall, the document first establishes the clinical and technical context, then
presents the contributions and the implemented two-phase system together with its
empirical evaluation.